\subsection{Formatting CSS inside Java}\label{sec:FormattingCSSInsideJava}
In most cases, CSS\index{CSS}\index{Cascading Style Sheets|see {CSS}} \autocite{W3C:CSS} settings should be provided through a resource file\index{resource file}, because there they are easier to modify if necessary.

CSS\index{CSS} is not only used with HTML\index{HTML} in the context with browsers, but also for GUIs created with JavaFX\index{JavaFX} \autocite{JavaFX}.

Basically, a \bdq{Cascading Style Sheet}\index{CSS} consists from a set of rules. Each rule begins with a selector or list of selectors, followed by one or more declarations, wrapped in curly braces. A declaration consists of a property and a value, separated by a colon. The declaration is terminated by a semicolon.

This looks like this:
\begin{lstlisting}[language={}]
<selector> <selector> {
    <property>:<value>;
    <property>:<value>;
}

<selector> <selector> {
    <property>:<value>;
    <property>:<value>;
}
\end{lstlisting}

Or, when embedded into a Java source file:
\begin{lstlisting}
final var cssString = """
````<selector> <selector> {
````    <property>:<value>;
````    <property>:<value>;
````}

````<selector> <selector> {
````    <property>:<value>;
````    <property>:<value>;
````}\
````""";
\end{lstlisting}

Each rule starts on its own line. The list of selectors is terminated by the openening curly brace and a linefeed. The declarations will be indented with four spaces, terminated by the semicolon and a linefeed again. The closing curly braces is not indented and follows on a line of its own.

The rules will be separated by an empty line from each other. 

%==============================================================================
% End of file!!
%==============================================================================
